\documentclass[a4paper,11pt]{article}
\usepackage[T2A]{fontenc}
\usepackage[utf8]{inputenc}
\usepackage[russian]{babel}
\usepackage{amsmath,amssymb,mathtools}
\usepackage{cmbright}
\renewcommand{\familydefault}{\sfdefault}
\usepackage{icomma}
\usepackage[a4paper,left=16mm,right=16mm,top=12mm,bottom=11mm]{geometry}
\usepackage{microtype}
\usepackage{multicol}
\usepackage{xcolor}
\usepackage{hyperref}
\definecolor{accent}{HTML}{15747A}
\definecolor{darkaccent}{HTML}{164C55}
\definecolor{textgray}{HTML}{283238}
\definecolor{muted}{HTML}{65747A}
\definecolor{rulegray}{HTML}{CADADC}
\color{textgray}
\hypersetup{hidelinks}
\setlength{\parindent}{0pt}
\setlength{\parskip}{0pt}
\setlength{\columnsep}{8mm}
\setlength{\abovedisplayskip}{1mm}
\setlength{\belowdisplayskip}{1mm}
\emergencystretch=2em
\newcommand{\sect}[1]{\par\vspace{2.2mm}{\normalsize\bfseries\color{darkaccent}#1}\par\vspace{0.3mm}{\color{rulegray}\rule{\linewidth}{0.45pt}}\par\vspace{0.8mm}}
\newcommand{\fml}[1]{\par\noindent{\normalsize\(\displaystyle #1\)}\par\vspace{2.8mm}}
\begin{document}
\thispagestyle{empty}

\hypersetup{pdftitle={Главные формулы. Анастасия Павлова},pdfauthor={Лёвин Артём Александрович}}
\begin{center}
{\small\bfseries\color{accent}УЧЕБНАЯ ПАМЯТКА \quad\textbullet\quad ОСНОВНЫЕ ФОРМУЛЫ}\par
\vspace{1mm}
{\fontsize{21}{25}\selectfont\bfseries\color{darkaccent} Степени, движение и совместная работа\par}
\vspace{0.8mm}
{\small\color{muted}Анастасия Павлова\quad\textbullet\quad ЕГЭ · математика\quad\textbullet\quad 12.10.2026}
\end{center}
\vspace{2mm}
\noindent{\color{accent}\rule{\linewidth}{0.9pt}}
\vspace{1.3mm}
{\bfseries\color{darkaccent}Главные принципы решения}\par
\vspace{1mm}
\begin{multicols}{2}
\raggedcolumns
\small
\noindent\textbullet\quad Для текстовой задачи вводите одну неизвестную и выражайте остальные величины через неё.\par\vspace{0.75mm}
\noindent\textbullet\quad Для совместной работы складывайте производительности, время не складывайте.\par\vspace{0.75mm}
\columnbreak
\noindent\textbullet\quad Если уравнение рациональное, сначала задайте ограничения на знаменатели.\par\vspace{0.75mm}
\noindent\textbullet\quad После решения квадратного уравнения проверяйте физический смысл корней.\par\vspace{0.75mm}
\end{multicols}
\vspace{1.1mm}
\noindent{\color{accent}\rule{\linewidth}{0.5pt}}
\vspace{0.6mm}
\begin{multicols}{2}
\raggedcolumns
\normalsize
\sect{Степени}
\fml{a^m a^n=a^{m+n}}
\fml{(a^m)^n=a^{mn}}
\fml{a^{m/n}=\sqrt[n]{a^m}\;(a>0)}
\sect{Движение}
\fml{S=vt,\quad t=\frac Sv}
\fml{v_{\text{сбл}}=v_1+v_2}
\fml{v_{\text{ср}}=\frac{S_{\text{общ}}}{t_{\text{общ}}}}
\columnbreak
\sect{Совместная работа}
\fml{p=\frac1t\quad(A=1)}
\fml{\frac1T=\frac1{t_1}+\frac1{t_2}}
\fml{p_{\text{общ}}=p_1+p_2+p_3}
\sect{Квадратные уравнения}
\fml{D=b^2-4ac}
\fml{x_{1,2}=\frac{-b\pm\sqrt D}{2a}}
\fml{a\ne0}
\end{multicols}
\vfill
{\color{accent}\rule{\linewidth}{0.8pt}}\par\vspace{1.5mm}
{\small\color{muted}Лёвин Артём Александрович эксклюзивно для Анастасии Павловой}
\end{document}